\documentclass{beamer}
\usetheme{metropolis}

\title{LLM Limitations and Evaluation Rigor in Finance}
\author{Juan F. Imbet}
\institute{Paris Dauphine -- PSL University}
\date{}

\begin{document}

\begin{frame}
  \titlepage
\end{frame}

\section{Calibration and Overconfidence}

\begin{frame}{Calibration and Overconfidence}
  \begin{itemize}
    \item Calibration measures whether a model's stated confidence matches its empirical accuracy
    \item Reliability diagrams and Expected Calibration Error (ECE) quantify miscalibration
    \item LLMs in finance frequently exhibit overconfidence in probabilistic forecasts
    \item Post-hoc techniques such as temperature scaling can improve calibration
  \end{itemize}
\end{frame}

\section{Temporal Leakage and Look-Ahead Bias}

\begin{frame}{Temporal Leakage and Look-Ahead Bias}
  \begin{itemize}
    \item Pretrained LLMs may have memorized financial events that appear in test sets
    \item Point-in-time data discipline is essential to prevent inadvertent leakage
    \item Anonymization of company names and dates reduces memorization risk
    \item Contamination-resistant splits use strict temporal cutoffs and holdout anonymization
  \end{itemize}
\end{frame}

\section{Stock Movement Prediction: A Cautionary Tale}

\begin{frame}{Stock Movement Prediction: A Cautionary Tale}
  \begin{itemize}
    \item Efficient market theory predicts that publicly available news should yield no edge
    \item Empirical studies report 51--65\% accuracy on short-horizon stock movement tasks
    \item ARIMA and LSTM baselines often match or exceed LLM accuracy at lower cost
    \item Reported gains frequently vanish after correcting for look-ahead bias and transaction costs
  \end{itemize}
\end{frame}

\section{Evaluation Beyond Classification Accuracy}

\begin{frame}{Evaluation Beyond Classification Accuracy}
  \begin{itemize}
    \item Sharpe ratio, alpha generation, and information ratio measure economic value
    \item Walk-forward and expanding-window validation mimic live trading conditions
    \item Transaction costs, market impact, and capacity constraints erode paper returns
    \item Maximum drawdown and tail-risk metrics capture downside beyond mean returns
  \end{itemize}
\end{frame}

\section{Hallucination in Financial Contexts}

\begin{frame}{Hallucination in Financial Contexts}
  \begin{itemize}
    \item LLMs fabricate citations, invent numerical figures, and hallucinate company details
    \item Retrieval-Augmented Generation (RAG) grounds outputs in verified source documents
    \item Retrieval verification and tool use add an additional layer of factual grounding
    \item Domain-specific hallucination benchmarks are needed to measure and compare risk
  \end{itemize}
\end{frame}

\end{document}
